\documentclass[aps,prx,reprint,superscriptaddress,floatfix,nofootinbib,longbibliography,amsmath,amssymb]{revtex4-2}
\usepackage[T1]{fontenc}
\IfFileExists{lmodern.sty}{\usepackage{lmodern}}{}
\usepackage{amsthm,mathtools}
\usepackage{enumitem}
\usepackage{booktabs,array}
\usepackage{xcolor}
\usepackage{tenkz}
\usetikzlibrary{quantikz2}
\usetikzlibrary{arrows.meta,calc,patterns,positioning,decorations.pathreplacing}
\usepackage{graphicx,adjustbox}
\usepackage[colorlinks=true,linkcolor=blue!50!black,citecolor=blue!50!black,urlcolor=blue!50!black]{hyperref}

% ---------- notation ----------
\newcommand{\C}{\mathbb C}
\newcommand{\R}{\mathbb R}
\newcommand{\Z}{\mathbb Z}
\newcommand{\E}{\mathbb E}
\newcommand{\id}{\mathrm{I}}
\newcommand{\Tr}{\operatorname{Tr}}
\newcommand{\dist}{\operatorname{dist}}
\newcommand{\supp}{\operatorname{supp}}
\newcommand{\rank}{\operatorname{rank}}
\newcommand{\diam}{\operatorname{diam}}
\newcommand{\norm}[1]{\lVert #1\rVert}
\newcommand{\abs}[1]{\lvert #1\rvert}
\providecommand{\ket}[1]{|#1\rangle}\renewcommand{\ket}[1]{|#1\rangle}
\providecommand{\bra}[1]{\langle #1|}\renewcommand{\bra}[1]{\langle #1|}
\newcommand{\cH}{\mathcal H}
\newcommand{\cB}{\mathcal B}
\newcommand{\cE}{\mathcal E}
\newcommand{\cI}{\mathcal I}
\newcommand{\cR}{\mathcal R}
\newcommand{\cV}{\mathcal V}
\newcommand{\cT}{\mathcal T}
\newcommand{\cA}{\mathcal A}
\newcommand{\cF}{\mathcal F}
\newcommand{\cS}{\mathcal S}
\newcommand{\cQ}{\mathcal Q}
\newcommand{\cK}{\mathcal K}
\newcommand{\Om}{\Omega}
\newcommand{\Lam}{\Lambda}
\newcommand{\e}{\mathrm e}

\theoremstyle{plain}
\newtheorem{theorem}{Theorem}[section]
\newtheorem{proposition}[theorem]{Proposition}
\newtheorem{lemma}[theorem]{Lemma}
\newtheorem{corollary}[theorem]{Corollary}
\theoremstyle{definition}
\newtheorem{definition}[theorem]{Definition}
\newtheorem{remark}[theorem]{Remark}
\newtheorem{construction}[theorem]{Construction}

% ---------- parts (REVTeX has no \part) ----------
\newcounter{notepart}
\renewcommand{\thenotepart}{\Roman{notepart}}
\newcommand{\notepart}[2]{\refstepcounter{notepart}\label{#2}%
  \par\medskip\noindent{\centering\large\bfseries Part~\thenotepart\\[2pt]#1\par}\medskip}

% ---------- the "written twice" device ----------
% Every mathematical unit of this note appears twice: first in words
% (what it says, why it is true, what it is for), then precisely.
\newenvironment{inwords}{%
  \par\medskip\noindent
  \begin{list}{}{\leftmargin=0.8em\rightmargin=0.8em}
  \item\relax
  \noindent{\sffamily\bfseries\color{blue!40!black}In words.}\ \sffamily\small}{%
  \end{list}\par\medskip}
\newcommand{\precisely}{\par\noindent{\bfseries Precisely.}\ }

\begin{document}

\title{A one-dimensional area law and a matrix product state
approximation from a global spectral gap: the collar-free route}
\author{Prepared for Sirui Lu}
\affiliation{Max Planck Institute of Quantum Optics, Garching}
\date{October 16, 2026}

\begin{abstract}
\noindent
We prove an entanglement area law and a polynomial matrix product state
approximation for the unique gapped ground state of any finite-range
one-dimensional Hamiltonian, specialising the two-dimensional arguments
of \emph{A two-dimensional area law from a global spectral gap} \cite{AL}
and \emph{Polynomial PEPS approximation of gapped square-grid ground
states} \cite{PP} to a spin chain. Two structural simplifications occur on
a line. First, the area-law proof collapses to three elementary steps: the
marginal tail bound and the buffered conditional-entropy discount of
\cite[Section~2]{AL}, combined with the Araki--Lieb inequality, directly
yield $S(X)\le C(q,R)\,(1+J/\Delta)$ for every interval $X$. The
reason is geometric: a shell around an interval is crossed by a bounded
number of interactions, so the conditional-entropy discount costs a
constant rather than a boundary-scaled error, and the
buffer whose entropy must be controlled contains only a constant number of
sites. The sub-volume recursion, positive quasi-local constraints, replica
scanner, amplification, and two-family geometry of \cite{AL} drop out
entirely. Second, the tensor-network construction reduces from a dyadic
hierarchy to a linear left-to-right \emph{sweep}. The encoded holes,
adaptive patch constraints, angular information estimates, and row--column
routing of \cite{PP} disappear; what remains are Uhlmann-type birth,
death, and exchange moves, an information reset via a two-sided Gaussian
spectral filter, dimension-free Gaussian density compression, and column
extraction. The output is a matrix product state of bond dimension
$CL^{c}$ within vector error $L^{-1}$ of the ground state. We also record
that in one dimension the marginal tail bound makes the matrix product
state theorem a two-line corollary of the area law by sequential Schmidt
truncation. Every statement in this note is written twice: first in words,
then precisely. The area-law proof uses none of the techniques of Hastings
or of Arad--Kitaev--Landau--Vazirani; the tensor-network protocol uses a
Lieb--Robinson bound in a single lemma, and we indicate how to avoid it at
the price of a quasi-polynomial bond dimension.
\end{abstract}
\maketitle



\section{Overview}

\begin{inwords}
In a gapped quantum lattice system, entanglement across a region is
generated along its boundary. In two dimensions, a region of linear size
$r$ has a boundary of $O(r)$ sites: every variational estimate across the
boundary incurs an $O(r)$ error, and any buffer introduced to isolate the
region is an annulus whose own entanglement is unknown. The multi-scale
replica, scanner, and amplification apparatus of \cite{AL} exists to
break this circularity. In one dimension, the boundary of an interval is at
most two points. Any variational operation across the boundary disturbs a
bounded number of interactions and costs $O(1)$ energy, while a buffer of
constant width has, trivially, constant entropy. The circularity never arises: a marginal tail bound,
a buffered conditional-entropy discount, and the Araki--Lieb inequality
establish the area law in full. This is the content of
Part~\ref{part:arealaw}.

Constructing a global tensor-network approximation requires more than
entropy bounds alone (the papers explain why: an entropy bound on each cut
does not give a global vector approximation).
The two-dimensional construction of \cite{PP} therefore distributes the
ground state among ``parties'' attached to dyadic squares via local
changes of ownership controlled by mutual information, and then compresses
the resulting distributed density operator into a tensor network. On a
line, a dyadic hierarchy is the wrong geometry: a binary tree of intervals
has cutwidth $\log L$ across the middle of the chain, which turns a
polynomial link dimension into a quasi-polynomial bond dimension. The
right geometry is a left-to-right sweep with a moving ``remainder'' party,
which keeps the cutwidth bounded. The only topological obstruction during
the sweep is dissolving a narrow remainder band wedged between two
distinct neighbouring owners; this is resolved by the same sequence of
Uhlmann birth, death, and two-sheet exchange moves used along block edges
in \cite{PP}. This is Part~\ref{part:mps}. We also show that, in one
dimension, the marginal tail bound of Part~\ref{part:arealaw} makes the
matrix product state theorem an immediate corollary of the area law by
sequential Schmidt truncation; the sweep protocol is the one-dimensional
counterpart of the two-dimensional construction, and it is the mechanism
that generalises to higher dimensions.

Table~\ref{tab:map} summarises how the two-dimensional architecture
specialises to a line. Readers interested solely in the one-dimensional
area law need only read Sections~\ref{sec:setting}--\ref{sec:closing}
(about four pages).
\end{inwords}

\begin{table*}[t]
\centering\small
\begin{tabular}{@{}p{0.33\textwidth}p{0.63\textwidth}@{}}
\toprule
Two-dimensional ingredient & Specialisation to one dimension\\
\midrule
Marginal tail bound (\cite{AL}, Lemma 2.1) & Kept verbatim; the cut parameter $\mathcal B$ is a constant (Section~\ref{sec:tail}).\\
Buffered conditional-entropy discount (\cite{AL}, Lemma 2.2) & Kept; the buffer has \emph{constant} width and the error is a constant (Section~\ref{sec:buffer}).\\
Sub-volume recursion, exponent $e_0<1$ & Unnecessary: the Araki--Lieb inequality closes the argument directly (Section~\ref{sec:closing}). Sketched for comparison in Section~\ref{sec:longroute}.\\
Positive quasi-local constraints, truncation & Unnecessary for the area law; replaced by the Gaussian filter reset in the MPS protocol.\\
Conditional estimates, replicas, transport, comparators, scanner & Unnecessary.\\
Amplification by random constraint products & Unnecessary; Section~\ref{sec:longroute} explains how it would apply on a line.\\
Two families of regions & Trivial on a half-line; Section~\ref{sec:longroute}.\\
\midrule
Coarse collar and angular information (\cite{PP}, Section 3) & Replaced by a direct consequence of the area law and information reset (Lemma~\ref{lem:logsep}).\\
Information reset by Gaussian filter (\cite{PP}, Lemma 3.3) & Kept with one-dimensional Lieb--Robinson locality (Section~\ref{sec:reset}).\\
Adaptive patch constraints, encoded holes (\cite{PP}, Sections 4, 6.1) & Unnecessary: no three owners meet on a line, and logarithmic blocks have polynomial Hilbert-space dimension.\\
Births, deaths, exchanges (\cite{PP}, Lemmas 6.6--6.8) & Kept verbatim (Section~\ref{sec:moves}).\\
Dyadic hierarchy with edge bands and point treatments (\cite{PP}, Section 7) & Replaced by a left-to-right linear sweep (Section~\ref{sec:sweep}).\\
Density compression, column extraction (\cite{PP}, Sections 5, 8) & Kept verbatim (Section~\ref{sec:compression}).\\
Row--column routing with constant congestion & Replaced by direct placement along the chain (Section~\ref{sec:routing}).\\
\bottomrule
\end{tabular}
\caption{Correspondence between the two-dimensional constructions of \cite{AL,PP} and their one-dimensional specialisation in this note: what survives, what simplifies, and what drops out.}
\label{tab:map}
\end{table*}

\precisely
Throughout, $q\ge 2$ and $R\ge 1$ are integers and $J,\Delta>0$. A finite
set $\Lam\subset\Z$ carries the induced path graph (with edges between
consecutive integers that both lie in $\Lam$) and graph distance $d_\Lam$.
Each site carries $\C^q$ and $\cH_\Lam=\bigotimes_{x\in\Lam}\C^q$.
An admissible Hamiltonian is
\begin{equation}\label{eq:ham}
 H=\sum_{\substack{\varnothing\ne X\subseteq\Lam\\ \diam_{d_\Lam}X\le R}}h_X,
 \qquad h_X=h_X^*,\qquad \norm{h_X}\le J,
\end{equation}
with one summand per support and zero summands allowed. The standing
spectral hypothesis is that $H$ has a unique normalised ground vector $\Om$
(up to phase), with ground energy $E_0$ and spectral gap $\Delta$:
\begin{equation}\label{eq:gap}
 H-E_0\id\ \ge\ \Delta\bigl(\id-\ket\Om\bra\Om\bigr).
\end{equation}
For $A\subseteq\Lam$, let $\rho_A=\Tr_{\Lam\setminus A}\ket\Om\bra\Om$ be
the reduced density operator and $S(A)=S_\Om(A)=-\Tr\rho_A\log\rho_A$ its
von Neumann entropy (with natural logarithms and $0\log 0=0$). For
disjoint $A,B\subseteq\Lam$ (writing $AB$ for $A\cup B$), the conditional
entropy and mutual information are $S(A|B)=S(AB)-S(B)$ and
$I(A:B)=S(A)+S(B)-S(AB)$. The edge boundary of $A$ in $\Lam$ is
$\partial_\Lam A=\{\{x,x+1\}\subseteq\Lam:\ \text{exactly one of }x,x+1\text{ lies in }A\}$.
Constants $C,c$ depend only on $q,R,J,\Delta$ unless stated otherwise and
may change from line to line; where a specific dependence is tracked, it
is stated explicitly. The two main theorems of this note are as follows.

\begin{theorem}[One-dimensional area law]\label{thm:A}
There is $C_{\mathrm{AL}}=C_{\mathrm{AL}}(q,R,J,\Delta)$, with
\[
 C_{\mathrm{AL}}\le C_0(q,R)\,\Bigl(1+\frac J\Delta\Bigr),
\]
such that for every finite $\Lam\subset\Z$, every admissible Hamiltonian
satisfying \eqref{eq:gap}, and every $A\subseteq\Lam$,
\[
 S_\Om(A)\le C_{\mathrm{AL}}\,\abs{\partial_\Lam A}.
\]
\end{theorem}

\begin{theorem}[One-dimensional matrix product state approximation]\label{thm:B}
Fix $q,R,J,\Delta$. There are $C,c>0$ depending only on these parameters
such that for every $L\ge2$ and every admissible Hamiltonian on
$\Lam=\{1,\dots,L\}$ satisfying \eqref{eq:gap}, there is a nonzero matrix
product state $\Phi$ on $\Lam$ with open boundary conditions and bond
dimension at most $CL^{c}$ such that
\[
 \min_{\theta\in\R}\Bigl\|\frac{\Phi}{\norm\Phi}-\e^{i\theta}\Om\Bigr\|\le L^{-1}.
\]
\end{theorem}

The constants never depend on $\Lam$, on $A$, or on the specific
Hamiltonian. No spectral gap for restricted sub-Hamiltonians, no
frustration-freeness, and no translation invariance are assumed.

\paragraph{Formalization.}
Theorems~\ref{thm:A} and~\ref{thm:B}, with explicit constants, are
formalized in Lean~4 in the TNLean modules \texttt{TNLean/MPS/AreaLaw/}
(blueprint chapter ``The one-dimensional area law and matrix product state
approximation''), stated directly on the chain of $L$ sites and, for
Theorem~\ref{thm:A}, on finite subsets of the integers. The counting
constants $N_R$, $\ell_R$, $\cB_R$ and the shell width $R$ are those printed
here. The buffered discount is run at the energy threshold $\Delta/2$ in
place of $\Delta/4$, the shell count omits the last entry of \eqref{eq:CE},
and the constant $1+2C_E/\Delta$ of \eqref{eq:discount} is replaced by
$4096\,\e\,\vartheta\,\cB_R/m+m\log 2$, which changes $C_{\mathrm{int}}$ and
$C_0$ accordingly; these deviations are recorded in
\nolinkurl{docs/paper-gaps/tnlean_one_dimensional_discount_constants.tex}.
Theorem~\ref{thm:A} is formalized for every finite $\Lam\subset\Z$; the
component reduction of Section~\ref{sec:setting} is carried out through the
vanishing entropy of the unions of components rather than the factorization
of $\Om$
(\nolinkurl{docs/paper-gaps/tnlean_one_dimensional_connected_domain_scope.tex});
Theorem~\ref{thm:B} is formalized both by the two-line route of
Section~\ref{sec:shortcut} and by the physical protocol of Part~\ref{part:mps}.
Proposition~\ref{prop:VC} is proved for every prescribed bond dimension $D$
and every orthonormal marginal eigenbasis ordered by nonincreasing
eigenvalue, including choices within repeated eigenspaces, with its
discarded-mass error bound before normalization. The threshold form used in
the two-line argument remains a separate result
(\nolinkurl{docs/paper-gaps/tnlean_one_dimensional_threshold_truncation.tex}).
The physical preparation is constructed on sufficiently large contiguous
chains, with half-open lattice bands and only the main registers as physical
outputs; auxiliary copies remain private at their corrected owners
(\nolinkurl{docs/paper-gaps/tnlean_one_dimensional_band_sweep_corrections.tex}).
Its logarithmic width constant is selected at information exponent $60$
and may be enlarged; the result is not asserted for every printed choice
at exponent $24$. Preparation and compression, with exact representations
at the finitely many shorter lengths, give a unit matrix product state of
polynomial bond dimension and phase error at most $1/L$ for every $L\ge2$.
The constants precede the length, Hamiltonian and ground vector, and only
the original full-system gap is assumed.

